\section{A successful extraction can still lose the fact}

Consider an Inline XBRL amount whose literal is \code{1,234}, whose admitted
numeric transform treats the comma as grouping, whose scale is $-3$, and whose
context and unit identify a particular dollar observation. The resulting amount
is $1234\times10^{-3}=1.234$. Scaling that normalized value again gives 0.001234
and is wrong. An adapter must retain the literal, transform, scale, context,
unit, nil status and original occurrence so that this error can be found.

The same principle applies to dates and dimensions. A lexical instant of
31 December remains distinct from the exclusive 1 January boundary used for
an interval. Publication and ingestion are separate timestamps. Two equal
numbers with different contexts remain separate occurrences. A nil fact is
unknown, not zero. Dropping a typed dimension or context reference yields a
preservation failure for that fact even if the resulting Markdown looks clear.

The inspected Arelle interfaces support native XBRL contexts, dimensions,
transformations and validation. Its internal object IDs do not replace stable
source identity. The inspected Docling source offers useful layout and table
proposals, but its examined XBRL graph path does not retain every required
fact/context/unit/typed-dimension field. That path needs a preservation sidecar
or must report unsupported conversion for affected records. These are source
audit findings, not measured accuracy rankings of the tools.

Each adapter therefore pins code, dependency, model and taxonomy versions,
licenses, input and output schemas, transformations and field-loss behavior.
Preparation preserves bytes and receipts before parsing, inventories the
package, attaches occurrences and contexts, records normalization, checks
required fields independently and submits candidates for review. Partial
conversion retains both successful units and the exact failures or unexplored
frontier. Conversion counts cannot accept a package or an accounting treatment.

The proposed extraction comparison uses the same retained files, roots, budgets
and taxonomy resources. Manual located source units and transparent native
extraction provide the reference; applicable pdfplumber/Arelle configurations,
plain Docling and explicit context/footnote repairs form the candidate ladder.
Transcription, context, table structure, mapping, source support, retention,
historical admission and cost are separate outcomes. Lost hashes or locators,
look-ahead, silent omissions and untracked downloads veto the affected exchange.
The current 31 development annotations cannot establish a heldout accuracy
ranking, and cell counts within one report are not independent observations.

Reuse also has distribution conditions. Arelle's Apache and Docling's MIT code
licenses do not by themselves clear every model, taxonomy or dataset. ACTUS
technical-specification reuse retains its share-alike obligations; QuantLib,
Z3 and ledger tools retain their own code and dependency terms. The Beancount
and hledger comparison needs an explicit distribution boundary for copyleft
components. License admission records actual versions, notices, models and
data, rather than treating a repository's headline license as universal.
